\documentclass[11pt]{article}
\usepackage[margin=1in]{geometry}
\usepackage{amsmath,amssymb,amsthm,booktabs}
\usepackage[hidelinks]{hyperref}
\newtheorem{theorem}{Theorem}
\newtheorem{proposition}[theorem]{Proposition}
\newtheorem{corollary}[theorem]{Corollary}
\newcommand{\mr}{\operatorname{mr}}
\newcommand{\Z}{\mathbb Z}
\newcommand{\F}{\mathbb F}
\title{A determinant obstruction to tree Coxeter quotients and fishnet scope}
\author{Internal research working proof}
\date{17 November 2026}
\begin{document}
\maketitle

The twenty-four sharp-four knots certified by seven-sheeted covers
are outside the sufficient fishnet family in Blair--Kjuchukova--Pfaff,
Proposition~1.4 \cite{bkp}.  The obstruction is an invariant of the
knot, so it excludes every possible fishnet presentation satisfying
that proposition.  Each such presentation would force the determinant
to have at least three prime factors counted with multiplicity.
The twenty-four determinants have at most two.

The argument applies more generally to reflection quotients whose
finite relation graph is a tree.  It does not exclude every fishnet
construction, every twisted-link theorem, or all prior resolutions of
these knots.  The group argument is a written proof; the finite
determinant computation is internally replayed, without formal
verification or external specialist review.  No global priority claim
is made.

\section{The orientation subgroup of a Coxeter presentation}

Let $G$ be a connected finite simple graph with vertices
$\{0,\ldots,r-1\}$.  Give each edge $e$ a finite integer label
$m_e\ge2$, and set
\[
 W(G,m)=\langle a_0,\ldots,a_{r-1}\mid
          a_v^2=1,\ (a_u a_v)^{m_{uv}}=1\text{ for }uv\in E(G)\rangle.
\]
There is \emph{no} relation on nonedges: their Coxeter exponent is
infinite.  Thus an omitted edge in this finite-relation graph does
not encode commuting generators.  Let $W^+$ be the kernel of the
homomorphism $W\to C_2$ sending every $a_v$ to the nontrivial element.

\begin{proposition}\label{prop:weighted}
The abelianization of $W^+$ has the integral presentation
\[
 H_1(W^+;\Z)=\Z^{r-1}/
 \langle m_{uv}(t_v-t_u):uv\in E(G)\rangle,
 \qquad t_0=0.
\]
Its finite order is
\[
 h(G,m)=\gcd_{T\text{ spanning tree of }G}\ \prod_{e\in E(T)}m_e.
\]
If $G$ itself is a tree, then in fact
\[
                         W^+\cong *_{e\in E(G)} C_{m_e}.
\]
\end{proposition}
\begin{proof}
Before the edge relations, the even subgroup of $*_{v=0}^{r-1}C_2$
is free on $t_v=a_0a_v$, $v\ne0$.  This follows directly by
Reidemeister--Schreier rewriting with transversal $\{1,a_0\}$.
Writing $t_0=1$ multiplicatively, an edge relation rewrites to
$(t_u^{-1}t_v)^{m_{uv}}$.  Its other transversal rewrite is
$(t_u t_v^{-1})^{m_{uv}}$, which is a conjugate of the inverse
of the first: conjugate $(t_v^{-1}t_u)^{m_{uv}}$ by $t_u$.
Hence only the first rewritten relation is needed for each edge.
Abelianization gives the displayed weighted incidence presentation.

Every maximal minor of the reduced incidence matrix is zero unless
its selected edges form a spanning tree.  For a spanning tree the
unweighted minor is $\pm1$, so the weighted minor is
$\pm\prod_{e\in E(T)}m_e$.  The greatest common divisor of all
maximal minors is the index of the full-rank relation lattice,
giving $h(G,m)$.

If $G$ is a rooted tree, orient its edges away from the root and
put $x_{uv}=t_u^{-1}t_v$.  These edge words are a free basis:
each $t_v$ is the ordered product of the $x_e$ along the root-to-$v$
path.  Thus the only defining relation on $x_e$ is $x_e^{m_e}=1$,
which proves the free-product assertion.
\end{proof}

\section{A necessary knot determinant condition}

Suppose a knot group admits an epimorphism onto $W(G,m)$ taking
meridians to reflections.  A reflection is conjugate to a standard
Coxeter generator, so it is an involution of odd sign.  The map
therefore factors through
\[
 \Pi_K=\pi_1(S^3\setminus K)/\langle\!\langle\mu^2\rangle\!\rangle
\]
and induces an epimorphism
\[
                       \pi_1(\Sigma_2K)\twoheadrightarrow W^+.
\]
Indeed the parity homomorphisms on source and target agree on every
meridian, so the even subgroup maps onto the even subgroup.

\begin{theorem}\label{thm:det}
Every such reflection quotient satisfies
\[
                         h(G,m)\mid\det K.
\]
For a tree on $r$ vertices this gives
\[
 \prod_{e\in E(G)}m_e\mid\det K,
 \qquad \Omega(\det K)\ge r-1,
\]
where $\Omega$ counts prime factors with multiplicity.  In particular,
all tree labels are odd, and a knot with prime determinant has no
tree reflection quotient of rank greater than two.
\end{theorem}
\begin{proof}
Abelianization preserves the epimorphism.  The finite group
$H_1(\Sigma_2K;\Z)$ has order $\det K$, while
$H_1(W^+;\Z)$ has order $h(G,m)$ by
Proposition~\ref{prop:weighted}.  The order of a finite quotient
divides the source order.  For a tree there are $r-1$ labels, each
at least two, so their product has at least $r-1$ prime factors
counted with multiplicity.  A knot determinant is odd; hence no
tree label can be even.
\end{proof}

The formula for general $G$ matters.  For a four-cycle with labels
$3,5,7,11$, the four spanning-tree products have greatest common
divisor one.  Such a cyclic relation graph supplies no determinant
divisibility obstruction at all.  The tree conclusion must not be
applied to arbitrary cyclic Coxeter diagrams.

\begin{corollary}\label{cor:weight}
If $3\nmid\det K$ and $K$ has a reflection quotient $W(G,m)$
on $r$ vertices, then
\[
                           \sum_{e\in E(G)}m_e\ge5(r-1).
\]
Consequently every BBK twisted diagram of a knot with
$3\nmid\det K$ has at least $5(b(K)-1)$ crossings.
\end{corollary}
\begin{proof}
The abelianization of $W$ is the binary vector space whose basis
is the set of connected components of the subgraph formed by the
odd-labelled edges.  Since a knot group is normally generated by
one meridian, this abelianization is cyclic, so that subgraph is
connected.  Choose a spanning tree $T_o$ in it.

By Theorem~\ref{thm:det}, $H_1(W^+;\F_3)=0$.  Reducing the
weighted incidence matrix modulo three shows that the edges with
$3\nmid m_e$ also form a connected graph.  Choose a spanning tree
$T_3$ there.  Put $k=|E(T_o)\cap E(T_3)|$.
An edge in the intersection has an odd label not divisible by
three, and therefore has weight at least five.  Each of the other
$r-1-k$ edges of $T_o$ has weight at least three, and each of the
other $r-1-k$ edges of $T_3$ has weight at least two.  Summing over
their union gives at least
$5k+3(r-1-k)+2(r-1-k)=5(r-1)$.

For a twisted diagram the edge weights are the absolute twist
numbers, their sum is its number of crossings, and BBK identify
$r$ with its knot's bridge number.
\end{proof}

All twenty-four determinants below are coprime to three, so any
twisted diagram for one of these bridge-four knots would have at
least fifteen crossings.  This statement does not assume that a
twisted diagram minimizes crossing number and does not exclude
a larger twisted diagram of a thirteen-crossing knot.

\section{The exact fishnet theorem excluded here}

In Definition~1.2 and Equation~(1.2) of \cite{bkp}, a regular fishnet
of width $2b$ has nonzero twist parameters $t_{i,j}$, and $d_j$ is
the positive greatest common divisor in column $j$.
Proposition~1.4 assumes
\[
                         d_2,d_4,\ldots,d_{2b-2}>1
\]
and proves $\mr(K)=b(K)=b$.  Its proof, Section~2.1, constructs
precisely the Coxeter quotient
\[
 \langle a_1,\ldots,a_b\mid a_i^2=1,
              (a_i a_{i+1})^{d_{2i}}=1\ (1\le i<b)\rangle.
\]
This finite-relation graph is a path.  Nonconsecutive pairs have
no imposed relation.

\begin{corollary}\label{cor:fishnet}
If a knot belongs to the sufficient family of Proposition~1.4 of
\cite{bkp}, then
\[
                     \Omega(\det K)\ge b(K)-1.
\]
None of the twenty-four seven-sheet sharp-four knots belongs to
that family, irrespective of the diagram or parameter choices.
\end{corollary}
\begin{proof}
The quotient just displayed and Theorem~\ref{thm:det} prove the
inequality.  For any of the twenty-four knots the independently
certified bridge number is four.  Therefore any presentation
satisfying Proposition~1.4 would have width eight and would force
$\Omega(\det K)\ge3$.  The exact calculations below give at most two.
\end{proof}

The determinant values are grouped only to shorten the table; each
literal PD is computed separately.
\begin{center}
\begin{tabular}{rlr}
\toprule
$\det K$ & Knot names & $\Omega$\\
\midrule
431 & $13a_{373},13a_{392}$ & 1\\
433 & $13a_{375},13a_{391}$ & 1\\
$19^2$ & $13a_{3264},13a_{3276}$ & 2\\
359 & $13a_{3266},13a_{3275}$ & 1\\
151 & $13n_{1005},13n_{1046}$ & 1\\
$13\cdot17$ & $13n_{1006},13n_{1047}$ & 2\\
59 & $13n_{1007},13n_{1048}$ & 1\\
137 & $13n_{1011},13n_{1043}$ & 1\\
211 & $13n_{1012},13n_{1044}$ & 1\\
$5\cdot17$ & $13n_{1013},13n_{1045}$ & 2\\
11 & $13n_{3842},13n_{3844},13n_{3855},13n_{3856}$ & 1\\
\bottomrule
\end{tabular}
\end{center}

\section{The nonunit plat gcd criterion is excluded as well}

For a $2n$-plat with braid $B$, let $R=\widehat\rho(B)$ be the
unreduced Burau representation evaluated at $t=-1$.  In the
conventions of Moriah's coauthor survey, let the $i$th row of $L$
be $e_{2i-1}^{\mathsf T}-e_{2i}^{\mathsf T}$ and the $i$th column
of $M$ be $e_{2i-1}+e_{2i}$.  Theorem~4.20 of \cite{moriah},
attributed there to Lustig--Moriah, says that if the greatest common
divisor $\alpha$ of all entries of $LRM$ differs from one, then
ordinary knot-group rank, exterior Heegaard genus and bridge number
all equal $n$.  Additional determinant congruences in that theorem
concern other conclusions and are not assumed here.

\begin{corollary}\label{cor:plat}
If a knot has a $2n$-plat for which $\alpha\ne1$, then
$\Omega(\det K)\ge n-1$.  Consequently none of the twenty-four
knots satisfies the nonunit plat gcd criterion of
\cite[Theorem~4.20]{moriah}, in any plat presentation.
\end{corollary}
\begin{proof}
Let $p$ be a prime dividing every entry of $LRM$; if all entries
are zero, choose any prime.  Assign arbitrary colours in $\F_p$
to the $n$ top cap pairs.  The map $M$ duplicates these colours,
$R$ propagates them through the braid by the Fox crossing rule,
and $L$ tests equality at the $n$ bottom cap pairs.  Hence the
Fox-colouring vector space is $\ker(LRM\bmod p)$.  Under the
divisibility assumption its dimension is $n$.

The Fox/core-group interpretation gives
\[
 \dim_{\F_p}\operatorname{Col}_p(K)
 =1+\dim_{\F_p}H_1(\Sigma_2K;\F_p).
\]
Thus $H_1(\Sigma_2K;\Z)$ has at least $n-1$ invariant factors
divisible by $p$, so $p^{n-1}$ divides $\det K$ and
$\Omega(\det K)\ge n-1$.

For the twenty-four knots, the cited criterion would give
$n=b(K)=4$, whereas the displayed determinants have
$\Omega\le2$.  This is impossible.
\end{proof}

This excludes the exact matrix gcd theorem. Its scope includes, and
extends beyond, the numerical criterion for generalized Montesinos knots.  It does
not substitute for an audit of every geometric family called
generalized Montesinos, or of results not requiring that gcd.

\section{Scope and verification}

The standard-library checker \texttt{verify\_fishnet\_scope.py}
reads the positive literal diagrams from the retained seven-sheet
inputs.  It constructs the integral Fox relation matrix
$2x_{\rm over}-x_{\rm in}-x_{\rm out}=0$, fixes one arc colour,
and computes all thirteen maximal minors by exact Bareiss
elimination.  The greatest common divisor is the order of the
presented finite abelian group $H_1(\Sigma_2K;\Z)$.
The Fox/core-group interpretation is the classical one already
sourced in the retained branched-cover proof \cite{core}.
The checker separately factors every determinant by exact trial
division and checks the tree and cyclic-graph controls above.

These checks rely on the retained diagram identity and bridge-upper
certificates for the named knots; they do not establish those
bindings anew.  They establish the determinant obstruction for
the supplied twenty-four knots, not an exhaustive Coxeter search.

Remark~1.3 of \cite{bkp} says every link can be drawn as a loose
fishnet.  Accordingly the conclusion is an exclusion from the
specific gcd family, not an assertion that these knots admit no
fishnet diagram.  Section~2.2 constructs other examples through
sublink direct sums, connected-sum constructions, and interweaving
fishnets.  Its displayed links and constructions require their own
hypotheses; an arbitrary Coxeter graph arising there need not be
a tree.  The determinant test applies whenever that finite-relation
graph is a tree, but does not by itself exclude every extension
described in that section.

The explicit constructions in that section introduce no additional
prime-knot cases beyond their input factors and the gcd family.
The direct-sublink constructions in Section~2.2.1 and the interweaving
construction in Section~2.2.2 retain at least two disjoint nonempty
sublinks, so their outputs have more than one component.  Observation~3
of Section~2.2.1 instead displays
$L_1\#T(2,t)\#L_2$.  For this to be a knot, $t$ must be odd and
both $L_i$ must be knots.  If the output is prime, all but at most
one of the three connected-sum factors must be unknots.  The
remaining factor is an input knot or the two-bridge torus knot.
This accounts for the displayed operations; it is not a classification
of every parameter generalization one might devise.

For the twisted-link theorem of Baader--Blair--Kjuchukova
\cite{bbk}, the Coxeter finite-relation graph is the weighted plane
dual of a twisted checkerboard surface.  Tree duals are covered
by Theorem~\ref{thm:det}; cyclic duals require the general weighted
incidence criterion.  The present note makes no assumption that
a twisted diagram realizes crossing number, and does not rule
out cyclic dual graphs by a crossing-count heuristic.

\begin{thebibliography}{9}
\bibitem{bkp} R. Blair, A. Kjuchukova and E. Pfaff,
Adding a suitable unknot to any link equates bridge number and
meridional rank, Math. Proc. Cambridge Philos. Soc. 181 (2026),
1339--1355. Definition~1.2, Equation~(1.2), Proposition~1.4
and its Section~2.1 proof; Section~2.2 separately.
\url{https://doi.org/10.1017/S030500412610214X}.
\bibitem{bbk} S. Baader, R. Blair and A. Kjuchukova,
Coxeter groups and meridional rank of links,
Math. Ann. 379 (2021), 1533--1551.
\url{https://doi.org/10.1007/s00208-020-02124-z}.
\bibitem{moriah} Y. Moriah, Heegaard splittings of knot exteriors,
Geometry \& Topology Monographs 12 (2007), 191--232,
Theorem~4.20 and its preceding definitions, reporting
Lustig--Moriah.  The coauthor survey's exact criterion is imported;
the full original proof is not independently audited here.
\url{https://doi.org/10.2140/gtm.2007.12.191}.
\bibitem{core} D. Silver, L. Traldi and S. Williams,
Core groups, J. Knot Theory Ramifications 33 (2024), 2450013,
introduction and Section~7, crediting Kelly and Wada.
\url{https://arxiv.org/abs/2311.02048v2}.
\end{thebibliography}
\end{document}
